\section{Background} \label{sec:background} 

This section describes how DAG-based agent workflows are mapped to an LLM serving system and how the system recovers an agent's KV state after a serving-worker failure.

\subsection{DAG-Based Agent Workflow Serving} \label{sec:background:serving} 

A multi-agent application coordinates multiple role-specialized agents to complete a user task. An agent may generate intermediate results, invoke external tools, delegate subtasks, and consume outputs produced by other agents. These dependencies form a directed acyclic graph (DAG), in which each node represents an agent invocation and each edge represents a data or control dependency. 

A \textbf{workflow} runtime maintains the DAG and releases an invocation after its dependencies have been satisfied. Forks allow multiple invocations to execute in parallel, while joins delay downstream invocations until all required predecessors have completed~\cite{lin2024parrot}. 
%
An \textbf{agent} can be invoked multiple times during one workflow. On each {turn}, its input may include an agent-specific instruction, previous model outputs, tool results, and outputs received from other agents. These inputs form an evolving agent history retained by the workflow runtime. 

Figure~\ref{fig:serving_model} shows how agent invocations are mapped to an LLM serving system. Ready invocations are submitted as inference requests and scheduled across model workers. A \textbf{worker} contains one model replica, or a distributed replica spanning a tensor-parallel~\cite{shoeybi2019megatron} group, and may concurrently serve agents from multiple workflows. 
%
Processing a request consists of prefill and decoding. During prefill, the model processes the input tokens and generates a key-value (KV) cache for each transformer layer. During decoding, newly generated tokens append entries to this cache. When a later invocation reuses a prefix of the same agent history, the corresponding KV entries allow the worker to process only the newly appended input rather than prefill the complete history again~\cite{gao2024cachedattention,zheng2024sglang}. Consequently, an active agent accumulates a KV history across its iterative turns. 

The KV cache is distributed according to the model-parallel~\cite{kwon2023vllm} configuration. Under tensor parallelism, each rank stores its local KV shards together with the metadata required by the attention implementation. Reusing a KV cache therefore requires compatibility in the model version, parallel configuration, numerical format, and KV layout. 


\begin{figure}[t]
  \centering
  \includegraphics[width=\columnwidth]{figures/serving_model.pdf}
  \caption{DAG-based agent workflow serving. Numbered agent invocations
    are requests; the runtime retains accepted histories and submits
    ready requests to a router for dispatch across workers. Each worker is one TP=2 replica.
    The illustrative MHA layer has eight Q/KV heads: ranks retain the
    same token positions for different heads. B1--B2 shows prefix reuse;
    hatched suffixes are new tokens. Local attention uses matching Q/KV
    heads, with output-projection contributions summed across ranks.
    Dashed arrows return accepted outputs. Bar lengths and placement
    are schematic; reuse requires matching, available KV.}
  \label{fig:serving_model}
  \Description{Two workflow DAGs sit above accepted token histories and
    tool results. Workflow 1 forks from A1 into B1 followed by B2, and
    into C1, then joins at A2. Workflow 2 forks from D1 into E1 and F1, then joins
    at D2. The central Router sends ready B2 and E1 requests to worker 0,
    and C1 and F1 to worker 1 along separate symmetric paths.
    Each worker contains two ranks. Rank 0 owns Q/KV heads 0 through 3
    and rank 1 owns heads 4 through 7 in an illustrative eight-head MHA
    layer. Each rank shows separate K and V bars for both assigned agents.
    For a given agent the two ranks cover the same token positions with
    different heads. Agent B's bars share a retained prefix and a hatched
    newly appended suffix. Dashed paths return outputs from both workers
    to the runtime. Internal computation is omitted from the drawing.}
\end{figure}


\subsection{Worker Failure and KV Recovery} \label{sec:background:recovery} 


We consider fail-stop failures that make a model worker, consisting of one model replica or a tensor-parallel group, unavailable. Such failures may originate from an accelerator, serving process, host, or communication fault~\cite{cao2026lumen,jayakody2026ghostserve}. The protected failure domain is deployment-specific and may correspond to a process, a distributed model replica, or an entire host. A recovery copy provides protection only if it remains reachable after the corresponding domain fails. For example, host-local DRAM or SSD may survive a process failure but does not protect against a host failure. 

A worker may hold KV states from multiple agents and workflows. Consequently, one worker failure can simultaneously interrupt several branches of the same workflow, as well as agents from different workflows. We treat all invocations whose KV states become unavailable under the same worker failure as \textit{one failure event}. 

We assume that the workflow DAG, accepted tokens, and completed tool results are retained outside the failed worker. Affected invocations can therefore be reissued to a compatible surviving replica with the same model version, parallel configuration, numerical format, and KV layout. We focus on model-worker failures because they remove large, worker-local KV state while leaving the workflow recoverable from its durable inputs. Failures of the workflow runtime or durable token store are orthogonal and can be handled using conventional replication mechanisms. Unlike these control and durable states, KV caches grow with agent histories, making full replication costly and creating the resource-allocation problem studied in this paper. We assume that at least one advertised recovery domain remains available. 

Without a surviving KV copy, the replacement replica must prefill each affected agent's retained history before generation can resume~\cite{strati2024dejavu}. Alternatively, the system can restore a complete KV checkpoint associated with an accepted token prefix and replay tokens accepted after that prefix~\cite{strati2024dejavu,cao2026lumen}. A complete checkpoint contains the historical KV tensors and attention metadata required by all model layers and model-parallel shards. A partial KV copy may accelerate restoration but cannot independently provide recovery coverage.
